\frontchapter{Bibliographie}
\begingroup
\setlength{\parindent}{0pt}
\setlength{\parskip}{.48em}
\newcommand{\bibentry}[1]{\par\hangindent=1.5em\hangafter=1 #1\par}
\origpage{540}{526}
\bibentry{AMUNTS (Katrin) et Al. \emph{Brain and Language}. Volume 89, 2004, 346 p.\ednote{原文如此。\emph{Brain and Language} 为期刊名；本条没有文章题名、完整作者或页码范围，尚不足以唯一识别所引论文。“346 p.”可能是引文定位，不宜当作全篇页数；保留待核。}}
\bibentry{Aristote. \emph{Catégories}. Flammarion, 2007.}
\bibentry{Aristote. \emph{De l’interprétation}. Éditions Les Échos du Maquis, 2014, 1(16 a), 3-8.}
\bibentry{Aristote. Métaphysique. Flammarion, 2007.}
\bibentry{Aristote. \emph{Premiers Analytiques}. Chapitre I, 24 B18-20. Flammarion, 2007.}
\bibentry{Aristote. \emph{Traité du ciel}. 295 B32, Flammarion, 2004.}
\bibentry{ARNAULD (Antoine) et NICOLE (Pierre). \emph{La Logique ou l’art de penser}. Paris, Presses Universitaires de France, 1965.}
\bibentry{ARRIVÉ (Michel). À la recherche de Ferdinand de Saussure. Paris, PUF, 2007, 20 p.}
\bibentry{AUROUX (Sylvain). \emph{La grammaire générale et les fondements philosophiques des classements de mots}. Langages, 1988, 85 p.}
\bibentry{AUSTIN (John Langshaw). \emph{Quand dire, c’est faire}. coll. « L’Ordre philosophique ». Paris, Éditions du Seuil, 1970, 37-45 p.}
\bibentry{BAKHTINE (Mikhaïl). \emph{Esthétique et théorie du roman}. Gallimard, 1978, 158-159 p.}
\bibentry{BARTHES (Roland). Éléments de sémiologie. Paris, Seuil, 1964, 131p.}
\bibentry{BARTHES (Roland). \emph{S/Z}. Paris, Seuil, 1976, 14-15 p.}
\bibentry{BENVENISTE (Émile). \emph{Bulletin de la Société Linguistique de Paris}. Tome LVI, Compte rendu des « Éléments ». Paris, Edouard Champion, 1960, 20-23 p.}
\bibentry{BENVENISTE (Émile). \emph{Problèmes de linguistique générale}. Chapitre I, Paris, Gallimard, 1966. 56 p.}
\bibentry{BLOOMFIELD (Leonard). \emph{Language}. Londres, George Allen \& Unwin Ltd, 1933.}
\bibentry{BOPP (Franz) et Al. \emph{Grammaire comparée}. Tome 5, BiblioBazaar, 2009.}
\bibentry{BOUHOURS (Dominique). \emph{Entretiens d’Ariste et d’Eugène}. Forgotten books, 2018.}
\bibentry{BRÉAL (Michel). \emph{Essai de sémantique : (Science des significations)}. Paris, Hachette, 1897, édition de 1924.}
\bibentry{BRÉAL (Michel). « Le langage et les nationalités. » dans \emph{Revue des deux mondes}, Paris, 1891, 615-639 p.}
\bibentry{BRÉAL (Michel). Les idées latentes du langage. Paris, Hachette, Leçon faite au Collège de France pour la réouverture du cours de grammaire comparée, le 7 décembre 1868, 296 p.}
\bibentry{BRÉAL (Michel). « Les lois intellectuelles du langage. Fragment de sémantique » dans \emph{Annuaire de l’association des études grecques en France}, tome 17, 133 p.}
\origpage{541}{527}
\bibentry{BRÉAL (Michel). Lettre à Angelo de Gubernatis. cité dans la revue Histoire, épistémologie, langage, 14 avril 1879, tome 3, 128p.\ednote{原文如此。此条把信件日期与后出的期刊卷次连排，未列转引论文的作者、题名、刊年；不能据此认定该期刊出版于1879年。暂保留各项，不猜补转引信息。}}
\bibentry{BRÉAL (Michel). « L’histoire des mots. » dans \emph{Revue des deux mondes}, Paris, 1er juillet 1887, 187-212 p.}
\bibentry{BRETON (André) et ÉLUARD (Paul). \emph{Dictionnaire abrégé du surréalisme}. Paris, José Corti, 1991.}
\bibentry{CARON (Jean). \emph{Précis de psycholinguistique}. PUF, Paris, 1997, 13 p.}
\bibentry{CHÂTEAU (Dominique). \emph{La théorie peircienne dans son cadre sémiotique : la question de l’icône}. Paris, L’Harmattan, 1997, 47 p.}
\bibentry{CHOMSKY (Noam). \emph{Aspects of the Theory of Syntax}. Cambridge, Massachusetts: MIT Press, 1965.}
\bibentry{CHOMSKY (Noam). \emph{La linguistique cartésienne}. Seuil, Paris, 1969.}
\bibentry{CHOMSKY (Noam). \emph{Le langage et la pensée}. Payot, 2009.}
\bibentry{CHOMSKY (Noam). \emph{Science, Mind, and Limits of Understanding}. Vatican, The Science and Faith Foundation (STOQ), Janvier 2014, 4 p.}
\bibentry{CHOMSKY (Noam). Un compte rendu du « Comportement verbal » de R. F. Skinner\ednote{原文如此。作者姓名的缩写应为 B. F. Skinner；同一总书目原书第531页列 SKINNER (Burrhus Frederic)。此处只订正姓名缩写，不把法译刊年改成英文原评刊年。} dans revue \emph{Langages}, N°16, 1969.}
\bibentry{CHOMSKY (Noam). « Three Factors in Language Design » dans \emph{Linguistic Inquiry}, Janvier 2005, 1-22 p.}
\bibentry{CHOMSKY (Noam). « Trois facteurs dans l’architecture du langage » dans \emph{Nouveaux cahiers de linguistique française}, 27, 2006, 1-32 p.}
\bibentry{CONDILLAC (Etienne d.). \emph{Traité des animaux}. Paris, Vrin, 2004.}
\bibentry{DAKHLIA (Jocelyne). \emph{Lingua Franca. Histoire d’une langue métisse en Méditerranée}. Actes Sud, 2008.}
\bibentry{DERRIDA (Jacques). \emph{Marges de la philosophie}. coll. « Critique », Paris, Éditions de Minuit, 1972, 382 p.}
\bibentry{DESCARTES (René). \emph{Méditation sixième}. Oeuvres de Descartes, Texte établi par Victor Cousin, Tome 1, Levrault, 1824, 322-350 p.}
\bibentry{DESCARTES (René). \emph{Méditations métaphysiques}. Première Méditation, Flammarion, 2009.}
\bibentry{DESCARTES (René). \emph{Oeuvres philosophiques}. Tome 3, Garnier, 885-887 p.}
\bibentry{DIK (Simon C.). \emph{The theory of Functional Grammar (Part 1: The structure of the clause)}. Edition Kees Hengeveld, Berlin, Mouton de Gruyter, 1997, 3 p.}
\bibentry{DORTIER (Jean-François). « Le débat Piaget/Chomsky » dans \emph{Sciences Humaines}, juillet 1999.}
\bibentry{DUCROT (Oswald). \emph{Le Dire et le dit}. Editions de Minuit, 1984.}
\bibentry{DUMARSAIS (César Chesneau). \emph{L’Encyclopédie}. Paris, 1ère édition, Tome 15, 1751, 16-34p.}
\bibentry{FENG (Yo-Lan). \emph{Histoire de la philosophie chinoise}. Princeton University Press, 1934, 206 p.}
\origpage{542}{528}
\bibentry{FERGUSON (Charles Albert). « Diglossia » dans \emph{Word}, vol.15, 1959, 435 p.\ednote{原文如此。该文刊于 Word 15(2)，页码为325–340；“435”不在该文页码范围内。此项据出版社提交的DOI书目元数据核对，不据此宣称重核全文引文。见\ecite{69}。}}
\bibentry{FINGER (Sarah). « Opération à cerveau “ouvert” » dans \emph{Journal du dimanche}, 13 mars 2003.}
\bibentry{FOUCAULT (Michel). \emph{La Grammaire générale de Port-Royal}. Didier, Larousse, 1967, 7-15 p.}
\bibentry{FOUCAULT (Michel). \emph{Les mots et les choses}. Paris, Gallimard, 1966, 70, 75 et 79 p.}
\bibentry{FREGE (Amending).\ednote{原文如此。Amending 不是此处作者的名字；应为 FREGE (Gottlob)。本书第四章正文与延伸书目均作 Gottlob Frege。所列选集、附录标题与版本信息尚待进一步核对，不能因姓名已订正便视作整条书目无误。} \emph{Appendice de Grundgesetze der Arithmetik}. Volume 2. Coll. « The Frege Reader ». Kessinger Publishing, LLC, 279 p.}
\bibentry{GADET (Françoise). \emph{La variation sociale en français}. Paris, Ophrys, 2003, 98 p.}
\bibentry{GARY-PRIEUR (Marie-Noëlle). \emph{De la grammaire à la linguistique}. L’étude de la phrase. Paris, Armand-Colin, 1985, 43 p.}
\bibentry{GAZDAR (Gerald). \emph{Pragmatics: Implicature, presupposition, and logical form}. New York, Academic Pr, 1979, 38 p.}
\bibentry{GLEASON (H.-A.). \emph{Introduction à la linguistique}. Larousse, 1969, 9-10 p.}
\bibentry{GOFFMAN (Erving). \emph{Les Rites d’interaction}. Éditions de Minuit, 1974, 9 p.}
\bibentry{GOTTLOB (Frege).\ednote{原文如此。姓与名倒置，应为 FREGE (Gottlob)，与第四章姓名及书目一致。} \emph{Que la science justifie le recours à une idéographie}. Coll. « Écrits logiques et philosophiques », Seuil, 1971, 63-64 p.}
\bibentry{GREGOIRE (Henri). \emph{Rapport sur la nécessité et les moyens d’anéantir les patois et d’universaliser l’usage de la langue française}. Séance du 16 prairial de l’an deuxième, Convention nationale, 4 juin 1794, 1-19 p.}
\bibentry{GREIMAS (A. J.). \emph{Du sens : Essais sémiotiques}. Paris, Seuil, 1970, 188p.}
\bibentry{GRICE (Paul). « Logique et conversation » dans revue \emph{Communication}, 1979, 57-72 p.}
\bibentry{GUILLAUME (Gustave). \emph{Faits de langue et faits de discours (1945-1946)}. Presses de l’Université Laval, 2009.}
\bibentry{HAGÈGE(Claude). \emph{Encyclopedia Universalis}, consulté le 21 juin 2019. URL : \url{http://www.universalis-edu.com/encyclopedie/andre-martinet/}}
\bibentry{HAHN (Carnap) et Al.\ednote{原文如此。Hahn 与 Carnap 是两位作者的姓，不是同一人的姓与名；应分别识别为 Hans Hahn 与 Rudolf Carnap。所列2010年卷由 Antonia Soulez 协调编订，另收 Neurath、Schlick、Waismann 等人的文本。见出版社记录\ecite{64}。} \emph{Manifeste du Cercle de Vienne et autres écrits}. Coll. « Bibliothèque des Textes Philosophiques », Vrin, 2010.}
\bibentry{HALLIDAY(Mickael). \emph{A Introduction to Functional Grammar}.\ednote{原文如此。标题中的不定冠词应为 An，而非 A。作者名 Mickael 应为 Michael，与本书第八章正文一致；标准姓名形式另见\ecite{65}。} London, Edward Arnold, 1985.}
\bibentry{HALLIDAY(Mickael). \emph{Learning how to mean: Explorations in the development of language}. London, Edward Arnold, 1975.}
\bibentry{HALLIDAY(Mickael). \emph{On Language and linguistics}. The Collected Works of M.A.K. Halliday Volume 3. London, Bloomsbury Academic \& Professional, 2006, 188p.}
\bibentry{HAMERS (Josiane) \& BLANC (Michel). \emph{Bilingualité et bilinguisme}. Bruxelles, Éditions Mardaga, 1983, 445 p.}
\bibentry{HENRY (Victor). \emph{Les Antinomies linguistiques}. Paris, Didier Erudition, 1896, 26 p.}
\origpage{543}{529}
\bibentry{HJELMSLEV (Louis Trolle). \emph{Essai d’une théorie des morphèmes, dans Essais linguistiques}. Paris, Editions de minuit, 1971, 161-173p.}
\bibentry{HJELMSLEV (Louis Trolle). « La Stratification Du Langage » dans \emph{Revue Word}, 1954, 163-165 p.}
\bibentry{HJELMSLEV (Louis Trolle). « La Structure morphologique » (1939) dans \emph{Essais linguistiques}. Paris, Editions de minuit, 1971, 133p.}
\bibentry{HJELMSLEV (Louis Trolle). \emph{Prolégomènes à une théorie du Langage}. Paris, Éd. de Minuit, 1971, chap. 22, 155p.}
\bibentry{HOVELACQUE (Abel). \emph{La linguistique}. Paris, Reinwald et Compagnie, 1876, 22-23 p.}
\bibentry{HUMBOLDT (William). \emph{Introduction à l’oeuvre sur le kavi}. éditions du Seuil, avril 1974.}
\bibentry{IGNACE (Kou P.-K.). \emph{Deux sophistes chinois : Houei Che et Kong-souen Long}. Paris, 1953, 163 p.}
\bibentry{JAKOBSON(Roman). \emph{Essais de linguistique générale}. Tome 1, Les fondements du langage. Paris, Éditions de Minuit, 1963.}
\bibentry{\emph{Journal de Chirurgie Viscérale}. Décembre 2010, Volume 147, Issue 6, 497-503 p.\ednote{原文如此。本条只有期刊、卷期与页码，缺作者和文章题名；不能把期刊名当作文章题名或作者。本版保留待核，不补造缺失字段。}}
\bibentry{KERBRAT-ORECCHIONI (Catherine). \emph{La conversation}. Paris, Seuil, 1997, 132 p.}
\bibentry{KRAFFT (Adolphe). \emph{Les Serments carolingiens de 842 à Strasbourg, en roman et tudesque}. E. Leroux, 1901.}
\bibentry{LABOV(William). \emph{A Life of Learning: Six People I Have Learned From}. American Council of Learned Societies, 2009, 1 p.}
\bibentry{LABOV(William). \emph{Sociolinguistique}. Éditions de Minuit, 1976, 259 p.}
\bibentry{LA METTRIE (Julien O. de) et Al. \emph{Oeuvres Philosophiques De La Mettrie}. Ulan Press, 2012.}
\bibentry{LAKOFF (Robin). « The logic of politeness » dans \emph{Papers from the ninth regional Meeting of the Chicago linguistic society}. T. Cedric Smith-Stark, \& A. Weiser (Eds.), 1973, 297-298 p.}
\bibentry{LEIBNIZ (Gottfried Wilhelm). \emph{Nouveaux Essais sur l’entendement humain}. Flammarion, 1921, 226-233p.}
\bibentry{LOCKE (John). \emph{Essai philosophique concernant l’entendement humain}. Paris, Vrin, trad. Coste, 1972, III, 2, § 1, 2.}
\bibentry{LÜDI (Georges) \& PY (Bernard). \emph{Etre bilingue}. Bern: Peter Lang. 146 p.}
\bibentry{MAILLARD (Michel). « Essai de typologie des substituts diaphoriques » dans \emph{Langue Française}, n° 21, février 1974, 55-71 p.}
\bibentry{MAROUZEAU (Jules). \emph{Lexique de la terminologie linguistique} (3e éd. : 1951). Paris, Geuthner, 1969, 71p.}
\bibentry{MARTINET (André). \emph{Connotations, poésie et culture, dans To Honor Roman Jakobson}. La Haye, Mouton, 1967, tome. 2, 1290 p.}
\bibentry{MARTINET (André). \emph{Éléments de linguistique générale}. Paris, Armand Colin, 1967.}
\origpage{544}{530}
\bibentry{MARTINET (André). \emph{Grammaire fonctionnelle du français}. Paris, Didier et Saint-Cloud, 1979, § 1, 32p.}
\bibentry{MARTINET (André). \emph{La linguistique synchronique}. Paris, PUF, 1965, 84p.}
\bibentry{MARTINET (André). \emph{La linguistique}. Paris, Denoël, 1969.}
\bibentry{MARTINET (André). \emph{Revue Roumaine de Linguistique}. Bucarest, Éditura Academiei Române, 1980, tome 25, 552p.}
\bibentry{MARTY (Claude) et MARTY(Robert). \emph{99 réponses sur la sémiotique}. Montpellier, CRDP, 1992, p.60.}
\bibentry{MEIGRET (Louis). \emph{Le tretté de grammere françoeze}. Len Pod, 2017.}
\bibentry{MONTAIGNE (Michel de). \emph{Essais. Livre II, Chapitre XII}. Flammarion, 1999.}
\bibentry{MORRIS (Charles). \emph{Signs, Language, and Behavior}. New York, G. Braziller, 1946, 219 p.}
\bibentry{MOUNIN (Georges) et Al. \emph{Dictionnaire de la linguistique}. PUF, 2004.}
\bibentry{OCHS (Elinor) \& SCHIEFFELIN (Bambi). \emph{Language acquisition and socialization: Three developmental stories}. Cambridge University Press, 1984, 289 p.}
\bibentry{PEIRCE (Charles Sanders). \emph{Collected Papers}. Cambridge Massachusetts, Harvard University Press, 1931.}
\bibentry{PEIRCE (Charles Sanders). Écrits sur le signe. Paris, Seuil, 1978.}
\bibentry{PEIRCE (Charles Sanders). « Elements of Logic (1903) » dans \emph{Collected Papers}. Cambridge Massachusetts, Harvard University Press, 1960, 247-249 p.}
\bibentry{PEIRCE (Charles Sanders). \emph{Lettres à Lady Welby}. New Haven. Conn, Whitlocks Inc., 1953, 422p.}
\bibentry{PEIRCE (Charles Sanders). \emph{Pragmatisme et Pragmaticisme}. Paris, Cerf, Coll. «Passages», 2002, 51p.}
\bibentry{PIAGET (Jean). \emph{Biologie et connaissance}. Paris, Gallimard, coll. L’avenir de la science, 1967, 20-25 p.}
\bibentry{PICOCHE (Jacqueline). \emph{Précis de lexicologie française}. Paris, Nathan, 1977, 104p.}
\bibentry{PINKER (Steven). \emph{The Language Instinct: How the Mind Creates Language}. New York, Harper Perennial Modern Classics, 1994.}
\bibentry{PLATON. « Cratyle » dans \emph{Platon : Œuvres complètes}. Éditions Garnier Flammarion, 2008, 67-69p.}
\bibentry{Porphyre. \emph{Isagoge}. Vrin, 2000.}
\bibentry{POTTIER (Bernard). \emph{Linguistique générale : théorie et description}. Paris, Klincksieck, 1974, 29p.}
\bibentry{RASTIER (François) Systématique des isotopies. dans Essais de sémiotique poétique. Antony, Larousse, éd par Julien Greimas 1972, 80-106p.}
\bibentry{RASTIER (François) \emph{Textes et sens}. CNRS-INALF et Paris-Sorbonne- Centre de Linguistique française, Paris, Didier Erudition, 1996, 44P.}
\bibentry{RASTIER (François), CAVAZZA (Marc) et ABEILLÉ (Anne). \emph{Sémantique pour l’analyse. De la linguistique à l’informatique}. Paris, Masson, Coll. Sciences cognitives, 1994, 38p.}
\bibentry{RASTIER (François). \emph{Sémantique interprétative}. Paris, PUF, 1987.}
\origpage{545}{531}
\bibentry{ROBERT (Estienne).\ednote{原文如此。本条不只是姓与名倒置：\emph{La précellence du langage françois} 的作者为 Henri Estienne，应按 ESTIENNE (Henri) 识别；下一条语法著作则属于 Robert Estienne。见\ecite{66}。} \emph{La précellence du langage françois}. Armand Colin, 1896.}
\bibentry{ROBERT (Estienne).\ednote{原文如此。此条姓与名倒置，应为 ESTIENNE (Robert)。2003年 Champion 校注本所据原著为1557年的 \emph{Traicté de la Grammaire Francoise}，参见\ecite{67}。} \emph{Traicté de la Grammaire Francoise}. Coll. « Textes de la Renaissance », Champion Honoré, 2003.}
\bibentry{ROMAINE (Suzanne). \emph{Bilingualism}. Oxford, Blackwell, 1995, 122 p.}
\bibentry{ROSIER (Irène). \emph{La Grammaire spéculative des modistes}. Presses Universitaires du Septentrion, 1983.}
\bibentry{ROUSSEAU (Jean) et THOUARD (Denis). \emph{Lettres édifiantes et curieuses sur la langue chinoise. Un débat philosophico-grammatical entre Wilhelm von Humboldt et Jean-Pierre Abel-Rémusat}. Lille, Presses Universitaires du Septentrion, 1999.}
\bibentry{RUSSELL (Bertrand). \emph{The Principles of Mathematics}. Cambridge University Press, 1903, 56 p.}
\bibentry{RUSSELL (Bertrand). \emph{The Principles of Mathematics}. Introduction, Chapitre III, 309-310 p.}
\bibentry{RUWET (Nicolas). « La grammaire générative. » dans \emph{Langages}, 1966, n°4, 6 p.}
\bibentry{RUWET (Nicolas). \emph{Introduction à la grammaire générative}. Paris, Plon, 1967.}
\bibentry{SAGET (Estelle). « Hugues Duffau : Le cerveau se répare tout seul » dans \emph{L’Express}, 2 octobre 2014.}
\bibentry{Salzburg Global Seminar. \emph{The Salzburg Statement for a Multilingual World}. décembre 2017.}
\bibentry{SALZMANN (Zdenek). Travaux du Cercle linguistique de Prague. John Benjamins Publishing Company, octobre 1999, Volume 1, 10 p.}
\bibentry{SAUSSURE (Ferdinand de). \emph{Cours de linguistique générale}. Paris, Payot, 1971.}
\bibentry{SAUSSURE (Ferdinand de). \emph{Cours de linguistique générale}. Paris, Payot, 1975, 38 p.}
\bibentry{SAUSSURE (Ferdinand de). \emph{Cours de linguistique générale}. Paris, Payot, 2005.}
\bibentry{SAUSSURE (Ferdinand de). \emph{Cours de linguistique générale}. Wiesbaden : Otto Harrassowitz, édition R. Engler, tome 2, 1967, 47p.}
\bibentry{Schleicher (August). \emph{De l’importance du langage pour l’histoire naturelle de l’homme}. Paris, Vrin, 1980, 80 p.}
\bibentry{SCHLEICHER (August). \emph{Fabel in indogermanischer Ursprache}. Berlin, Dümmler, 1868, 206-208 p.}
\bibentry{SCHLEICHER (August). \emph{La théorie de Darwin et la science du langage}. Paris, Vrin, 1980.}
\bibentry{SKINNER (Burrhus Frederic). \emph{Verbal Behavior}. Copley Publishing Group, 1991.}
\bibentry{Société de Linguistique de Paris. \emph{Bulletin de la Société de linguistique de Paris}. Volume 1, Article 2, 1871, 3 p.}
\bibentry{TESNIÈRE (Lucien). Éléments de syntaxe structurale. Paris, Klincksieck, 1959.}
\bibentry{THOMAS (Antoine). \emph{La sémantique et les lois intellectuelles du langage in Essais de philologie française}. Paris, Librairie Émile Bouillon, 1897, ch. XIX, 178p.}
\bibentry{THOUARD (Denis). « Humboldt, Abel-Remusat et le chinois : Du mystère au savoir » dans \emph{Revue Texto}, Université Charles-de-Gaulle Lille 3, 2001.}
\bibentry{TORT (Patrick). Évolutionnisme et linguistique. Paris, Vrin, 1980, 59-78 p.}
\origpage{546}{532}
\bibentry{TROUBETZKOY(N.-S.). \emph{Principes de phonologie}. Paris, Klincksieck, 1949, 81p.}
\bibentry{VAUGELAS (Claude F. de). \emph{Remarques sur la langue française. Utiles à ceux qui veulent bien parler et bien écrire}. Coll. « Champ libre », Ivréa, 1981.}
\bibentry{VOLOSHINOV (Valentin). \emph{Marxism and the Philosophy of Language}. 1973, Harvard University Press, 165 p.}
\bibentry{VON HUMBOLDT (Humboldt).\ednote{原文如此。括号内重复了姓氏，未给名字；本条所指为 Wilhelm von Humboldt，与本书第五章及原书第531页所列通信集中的姓名一致。} \emph{De l’origine des formes grammaticales et de leur influence sur le développement des idées}. Hachette Livre BNF, 2016.}
\bibentry{VON SCHLEGEL (August W.). \emph{Observations sur la langue et la littérature provençales}. Ulan Press, 2012, 13-14 p.}
\bibentry{VON SCHLEGEL (Friedrich). \emph{Essai sur la langue et la philosophie des Indiens}. Paris, Parent-Desbarres, 1837.}
\bibentry{WAISMANN (Friedrich). \emph{Ludwig Wittgenstein and the Vienna Circle}. Rowman \& Littlefield Publishers, 1979, 47 p.}
\bibentry{WALTER (Henriette). \emph{Honni soit qui mal y pense}. Paris, Robert Laffont, février 2001.}
\bibentry{WATSON (John B.). \emph{Behaviorism (Revised edition)}. Chicago, University of Chicago Press, 1930.}
\bibentry{William Marçais. « La diglossie arabe. » dans \emph{L’Enseignement public}, tome CIV, n° 12, 1930-1931, 401-409 p.}
\bibentry{WITTGENSTEIN (Ludwig). \emph{Le Cahier bleu et le Cahier brun}. Coll. « Bibliothèque de Philosophie », Paris, Gallimard, 1996, 56 p.}
\bibentry{WITTGENSTEIN (Ludwig). \emph{Recherches philosophiques}. Gallimard, 2005, 39-40 p.}
\bibentry{WITTGENSTEIN (Ludwig). \emph{Tractatus logico-philosophicus}. Proposition 7, Gallimard, 2001.}
\bibentry{方勇. 庄子. 中华经典名著全本全注全译丛书，中华书局，2015.}
\bibentry{黄克剑. 公孙龙子. 中华经典名著全本全注全译丛书，中华书局，2018.}
\endgroup
